\section{Experimental Design}
\label{sec:newexperiments}
\paragraph{Three adaptation settings.}
The starting encoder is OpenCLIP ViT-B/32 pretrained on LAION-2B \citep{cherti2023reproducible}.
Its original adapters are linear.
The nonlinear replication uses separate 128-unit ReLU branches with zero output initialization.
The second encoder is OpenAI RN50 with linear adapters \citep{radford2021clip}.
Its architecture, pretraining, and feature dimension change together.
Both encoders remain frozen.
Each replication crosses Source, Supported, and three stratified draws with three rates and three seeds: 45 fits.
Rates are $10^{-4}$, $3\times10^{-4}$, and $10^{-3}$; seeds are 17, 29, and 43.
Primary comparisons use epoch ten and identical rates.

\paragraph{A separate preservation study.}
The AD study uses linear adapters and the same 1,200 training images for both encoders.
It crosses $\lambda\in\{.5,.8\}$ and $\beta\in\{1,4,16\}$ with the declared rates and seeds.
Source, U, A, D, AD, and nine randomized controls require 126 fits per encoder.
Base fits have ten epochs; controls stop at their matched AD epoch.
Epoch zero remains available.
The earlier 180-fit design separately tests directional interventions and retention.
Appendix~\ref{app:newmethods} gives both full grids, interpolation candidates, and selection rules.

\paragraph{Development selection.}
Development retrieval uses 900 images and 4,500 source captions, excluding calibration images.
Its original 100-image validation subset supplies relation accuracy.
Each candidate and seed selects the most accurate epoch that satisfies both retrieval constraints.
Each direction can lose at most one percentage point against frozen initialization.
Shared coefficients and a shared rate maximize mean selected accuracy across seeds.
Fixed tie rules favor earlier epochs, lower rates, and lower coefficients.

\paragraph{Search and caption discrimination.}
Image queries retrieve one of five source captions: I$\to$T.
Caption queries retrieve their owning image: T$\to$I.
R@1 counts queries whose highest-scoring candidate is relevant.
The clean e-ViL pool contains 1,000 test images and 5,000 source captions.
Both directions search the complete pool.
SugarCrepe++ requires both valid captions to outrank an incorrect caption \citep{dumpala2024sugarcrepepp}.
Secondary endpoints cover SugarCrepe, relation accuracy, and COCO retrieval.
Appendix~\ref{app:newmethods} specifies their pools \citep{hsieh2023sugarcrepe}.

\paragraph{A joint practical criterion.}
All three selected seeds must contain nonzero updates.
Against frozen initialization, both adjusted retrieval bounds must exceed $-1$ percentage point.
The adjusted SugarCrepe++ gain must exceed zero.
The same family must pass for both encoders.
The AD study adjusts jointly over 80 primary contrasts and separately over 24 factorial contrasts.
Both use 100,000 paired image-cluster resamples.
Original retention keeps separate 80-effect and 36-effect families; each replication keeps its own twelve-effect family.
Caption queries stay with their image during resampling.
Intervals condition on fitted seeds, assignment draws, and fixed galleries.
